% chapters/01/checklist.tex
\bagian{Daftar Periksa Kesiapan Belajar}

Bab ini bersifat pengantar, sehingga daftar periksa berikut bukan daftar
kompetensi teknis, melainkan daftar kesiapan. Periksa satu per satu sebelum
memasuki Bab~III; butir yang belum tercentang menunjukkan hal yang perlu Anda
selesaikan lebih dahulu.

\begin{checklistbox}
\begin{checklist}
  \item Saya dapat menjelaskan ruang lingkup mata kuliah ini dan mengetahui
        materi apa saja yang berada di luar cakupannya.
  \item Saya dapat menyebutkan keempat CPMK beserta artefak yang menjadi
        buktinya.
  \item Saya memahami bahwa penilaian pada mata kuliah ini berbasis kriteria,
        bukan perbandingan dengan teman sekelas.
  \item Saya mengetahui tiga syarat kelulusan: nilai akhir minimal 60,
        kehadiran minimal 75\%, dan penyerahan seluruh keluaran wajib.
  \item Saya sudah membaca peta bab dan mengetahui bab yang menopang setiap
        materi pada RPS.
  \item Saya memahami prinsip penyelarasan konstruktif: capaian, materi,
        aktivitas, dan asesmen menunjuk kemampuan yang sama.
  \item Saya dapat menjelaskan alasan Python dan Pygame dipelajari sebelum
        Unity.
  \item Saya sudah menyiapkan, atau mengetahui cara menyiapkan, Python 3,
        editor kode, dan Git pada komputer yang akan saya pakai.
  \item Saya sudah membaca lima kebiasaan kerja pada
        Bagian~\ref{sec:cara-belajar} dan memilih kebiasaan yang akan saya
        latih lebih dahulu.
  \item Saya memiliki rencana waktu belajar mingguan di luar jadwal kuliah.
  \item Saya mengetahui ke mana harus bertanya ketika menemui kesulitan:
        dosen pengampu, asisten praktikum, atau dokumentasi resmi.
\end{checklist}
\end{checklistbox}

\begin{catatabox}
Praktikum pertama pada Bab~III menuntut Python dan Pygame sudah terpasang.
Apabila butir kedelapan pada daftar di atas belum tercentang, selesaikan
sebelum pertemuan pertama agar waktu praktikum tidak habis untuk pemasangan.
\end{catatabox}
